\section{Acetic acid và giấm ăn}

\subsection{Cấu tạo và tính chất vật lí}

Acetic acid có công thức phân tử \ce{C2H4O2}, công thức cấu tạo thu gọn
\[
\ce{CH3-COOH}.
\]
Nhóm \ce{-COOH} là nhóm đặc trưng của acetic acid.

Acetic acid là chất lỏng không màu, có vị chua, mùi đặc trưng và tan tốt trong nước. Dung dịch acetic acid loãng là thành phần tạo vị chua chính của giấm ăn.

\subsection{Tính chất hóa học của acetic acid}

Acetic acid làm quỳ tím hóa đỏ và có các phản ứng đặc trưng của acid.

\begin{ghinho}{Một số phản ứng cần viết thành thạo}
\begin{tasks}(1)
\task $\ce{2CH3COOH + Mg -> (CH3COO)2Mg + H2 ^}$.
\task $\ce{CH3COOH + NaOH -> CH3COONa + H2O}$.
\task $\ce{2CH3COOH + CuO -> (CH3COO)2Cu + H2O}$.
\task $\ce{2CH3COOH + CaCO3 -> (CH3COO)2Ca + CO2 ^ + H2O}$.
\end{tasks}
\end{ghinho}

Acetic acid cũng cháy hoàn toàn:
\[
\ce{CH3COOH + 2O2 -> 2CO2 + 2H2O}.
\]

\begin{vidu}
Cho \SI{12.0}{\gram} acetic acid tác dụng hoàn toàn với \ce{NaOH} dư. Tính khối lượng \ce{CH3COONa} thu được.
\end{vidu}
\begin{loigiai}
\[
n_{\ce{CH3COOH}}=\frac{12.0}{60}=\SI{0.200}{\mole}.
\]

\bangmol{5}
{\ce{CH3COOH} & + & \ce{NaOH} & \rightarrow & \ce{CH3COONa + H2O}}
{0.200 & & & & 0.200}

\[
m_{\ce{CH3COONa}}=0.200\times82=\SI{16.4}{\gram}.
\]

\textbf{Vậy,} thu được \SI{16.4}{\gram} \ce{CH3COONa}.
\end{loigiai}

\subsection{Giấm ăn và bài toán nồng độ}

Khi đề cho thể tích dung dịch, khối lượng riêng và nồng độ phần trăm, có thể đi theo chuỗi:
\[
V_{\text{dd}}\longrightarrow m_{\text{dd}}\longrightarrow m_{\text{chất tan}}.
\]

\begin{vidu}
Một chai giấm có thể tích \SI{1.00}{\liter}, khối lượng riêng \SI{1.045}{\gram\per\milli\liter}, chứa $4.0\%$ acetic acid theo khối lượng. Tính khối lượng acetic acid có trong chai.
\source{Phỏng theo dữ kiện đề tuyển sinh chuyên Hóa Lâm Đồng 2026.}
\end{vidu}
\begin{loigiai}
\[
m_{\text{dd}}=1000\times1.045=\SI{1045}{\gram}.
\]
\[
m_{\ce{CH3COOH}}=1045\times0.040=\SI{41.8}{\gram}.
\]

\textbf{Vậy,} chai giấm chứa \SI{41.8}{\gram} acetic acid.
\end{loigiai}

\subsection{Điều chế acetic acid từ ethylic alcohol}

Trong quá trình lên men giấm, ethylic alcohol bị oxi hóa bởi oxygen trong không khí:
\[
\ce{C2H5OH + O2 -> CH3COOH + H2O}.
\]

\begin{tuluyen}
Từ \SI{9.20}{\gram} ethylic alcohol, hiệu suất quá trình lên men thành acetic acid là $80\%$. Tính khối lượng acetic acid thu được.
\dapso{\SI{9.60}{\gram}.}
\begin{loigiaiBT}
\[
n_{\ce{C2H5OH}}=\frac{9.20}{46}=\SI{0.200}{\mole}.
\]

\bangmol{7}
{\ce{C2H5OH} & + & \ce{O2} & \rightarrow & \ce{CH3COOH} & + & \ce{H2O}}
{0.200 & & & & 0.200 & &}

\[
m_{\ce{CH3COOH},\text{lt}}=0.200\times60=\SI{12.0}{\gram},
\]
\[
m_{\ce{CH3COOH}}=12.0\times0.80=\SI{9.60}{\gram}.
\]

\textbf{Vậy,} thu được \SI{9.60}{\gram} acetic acid.
\end{loigiaiBT}
\end{tuluyen}
